\chapter{Conclusão}

\section{Síntese dos Resultados Obtidos}
Este trabalho propôs e avaliou um arcabouço metodológico sistemático focado no aprimoramento da detecção e classificação de falhas em cadeias de isoladores, utilizando técnicas de processamento de imagem aliadas a Redes Neurais Convolucionais (CNN). A metodologia estruturada permitiu avaliar o impacto de diversas técnicas de pré-processamento de forma isolada e combinada, bem como otimizar seus parâmetros fundamentais.

A análise empírica dos resultados demonstrou, de forma inequívoca, que o pré-processamento das imagens exerce um papel determinante no desempenho preditivo dos modelos em cenários de inspeção de infraestrutura elétrica. No entanto, constatou-se a não obviedade da escolha técnica: não existe um processamento universalmente superior, sendo a eficácia intrinsecamente condicionada às características do conjunto de dados e ao contexto de captura das anomalias.

Para o \textit{dataset} CPLID, cujas imagens apresentam fundos homogêneos e falhas estruturais demarcadas em materiais cerâmicos e de vidro, técnicas pautadas na simples amplificação do intervalo tonal, notadamente o \textit{Contrast Enhancement}, obtiveram resultados expressivos (alcançando acurácia de 95,35\% e \textit{F1-score} de 96,55\%). Paralelamente, confirmou-se que filtros de suavização espacial e desfoque (\textit{Gaussian Blur}, por exemplo) revelaram-se severamente comprometedores para a detecção morfológica das arestas do isolador. 

Em contrapartida, avaliando o \textit{dataset} DRNPW, composto sensivelmente por imagens processadas por drones e submersas em forte poluição visual e fundos ruidosos contendo áreas densas de vegetação ou ocupação geográfica nativa, o emprego exclusivo de abordagens simplórias retrocedeu. Nestas predições sensíveis, a arquitetura híbrida e combinatória revelou sua proficiência. A sucessão sinérgica pautada na aplicação concatenada entre realce (\textit{Contrast Enhancement}) e controle do espectro por equalização extensiva (\textit{Equalize Histogram}) destacou-se categoricamente na supressão da interferência do fundo da imagem. Essa abordagem resultou em um aumento substancial nas métricas de desempenho (acurácia e F1-score) quando comparada aos modelos treinados apenas com as imagens originais sem processamento.

Por fim, a averiguação da otimização algorítmica expôs plena demonstração do impacto de variáveis sutis no desempenho de redes profundas. Ao contrastar o rigor exaustivo da Busca em Grade (\textit{Grid Search}) com a flexibilidade da Busca Aleatória (\textit{Random Search}), observou-se uma notável concordância na identificação do ponto de ajuste ótimo, localizado em torno de fatores de realce entre 1,3 e 1,4. Esses resultados não apenas consolidam os ganhos quando comparados ao uso da imagem original, resultando na expressiva marca de 90\% de acurácia, mas também evidenciam a estabilidade do modelo nessa faixa. O mapeamento comprovou, contudo, que o uso de níveis externos excessivos de contraste distorce a imagem e introduz artefatos que prejudicam o aprendizado.

\section{Limitações e Desafios Encontrados}
A despeito das expressivas evidências a fator das implementações adotadas, o trabalho esbarrou em desafios inerentes à unificação entre a inspeção sistêmica aérea computacional e fluxos complexos modelares por aprendizagem profunda.

O principal desafio enfrentado foi a limitação na quantidade e na disponibilidade de dados. Como observado logo no início da pesquisa, há uma grande escassez de bancos de imagens públicos focados em redes de distribuição que contenham um número expressivo de componentes defeituosos com anotações rigorosas. Essa carência metodológica, comum na área de sensoriamento, muitas vezes obriga o uso de conjuntos de dados fortemente desbalanceados. Como consequência, o baixo volume de amostras da classe minoritária (isoladores com falhas) prejudica o treinamento eficiente do algoritmo, aumentando a tendência do modelo em focar na classe dominante e gerar falsos positivos a partir de elementos do contexto de fundo da imagem.

Outra limitação importante está no alto custo computacional exigido para encontrar os melhores hiperparâmetros. A metodologia proposta permite a criação iterativa de inúmeros \textit{pipelines}, combinando diversas técnicas de processamento em etapas sequenciais. Diante dessa enorme flexibilidade, tentar otimizar todos os parâmetros simultaneamente de forma exaustiva, testando estritamente cada combinação possível, resulta em um problema conhecido como explosão combinatória (fenômeno em que o número de combinações a serem testadas cresce de forma exponencial a cada nova variável adicionada). Na prática, devido a essa vasta quantidade de cenários gerados, o modelo exige um poder de processamento extremo que inviabiliza a avaliação empírica integral em um tempo hábil.

\section{Sugestões para Pesquisas Futuras}
Reconhecidas as limitações apontadas e inspirando-se nos apontamentos empíricos quantitativos atestados nas avaliações estatísticas dos conjuntos de testes perante os processamentos limpos das matrizes, elenca-se a seguir uma lista pragmática viável visando a continuidade científica natural das vertentes de evolução:

\begin{itemize}
    \item \textbf{Otimização Estocástica Multi-paramétrica Ampliada via Otimização Bayesiana:} Transcender os limites numéricos bidimensionais estáticos contidos nas buscas em grade puras valendo-se das propensões analíticas em matrizes preditivas orientadas como a inferência da Otimização Bayesiana, permitindo desbravar simultaneamente variações contínuas de múltiplos filtros iterativos dentro do \textit{pipeline} combinatório mitigando a referida penalidade de explosão processual.
    \item \textbf{Aumento Generativo da Densidade de Dados (GANs):} Contornar a supressão latente advinda da escassez nativa balanceadora utilizando redes geradoras antagonistas (Generative Adversarial Networks) aplicadas restritamente como mecanismo superamostral (emular sinteticamente defeitos microscópicos nas matrizes isoladoras perfeitas garantindo expansão representacional não enviesada e adensando os conjuntos esparsos).
    \item \textbf{Integração Cruzada às Topologias Sistêmicas Atuais:} Expandir o laço desenvolvido incorporando os processamentos empíricos isolados do crivo local em convoluções básicas de classificação rumo ao ambiente nativo exposto. Submeter e comparar abertamente a robustez paramétrica atingida pelo fluxo de contraste superior contrapondo a sua estabilidade nas predições de fronteira em topologias detentoras baseadas em segmentações da arquitetura multiescala tipo \textit{You Only Look Once} (YOLO nas suas especificações operacionais v8) providenciando o salto formal validador entre aprovação classificatória e a efetiva espacialidade analítica perimetral em campo.
    \item \textbf{Adequação Portável e Escalabilidade em Topologia de Borda:} Mapeamento temporal sistemático a fim de transpor os subprodutos consolidados rumo ao nível transiente de acoplamento físico aferindo ativamente o viés métrico temporal introduzido pelo complexo pipeline hiperparamétrico das imagens, confirmando que seu desempenho e fluidez se mantém aderente na portabilidade perante microsserviços inseridos embarcadamente nos hardwares visuais (Drones/VANTs), em consonância real validadora exigida ativamente pelos cronogramas da vigilância dinâmica diária executada.
\end{itemize}

\section{Considerações Finais}
Em suma, a principal contribuição deste trabalho reside no delineamento de uma metodologia estruturada para a inspeção automatizada de linhas de transmissão. Ao mapear sistematicamente o impacto gerado pelo processamento de imagens aéreas, o projeto comprovou a necessidade indispensável de se tratar os dados visuais brutos antes da etapa de modelagem.

O estudo demonstrou que não é suficiente investir apenas em arquiteturas complexas de redes neurais se as imagens de entrada sofrerem com ruídos e limitações de contraste que ocultam os defeitos. Portanto, as otimizações documentadas provam de forma objetiva que os modelos preditivos só atingem níveis de excelência quando o pré-processamento garante um ambiente visual mais nítido e destacado. É o tratamento adequado da imagem que fornece às redes convolucionais as condições propícias para aprender e identificar com precisão as anomalias e imperfeições em componentes mapeados em campo aberto.
